\documentclass[a4paper]{article}
% generated by Docutils <https://docutils.sourceforge.io/>
\usepackage{cmap} % fix search and cut-and-paste in Acrobat
\usepackage[T1]{fontenc}
\setcounter{secnumdepth}{0}
\usepackage{longtable,ltcaption,array}
\setlength{\extrarowheight}{2pt}
\newlength{\DUtablewidth} % internal use in tables
\newcommand{\DUcolumnwidth}[1]{\dimexpr#1\DUtablewidth-2\tabcolsep\relax}

%%% Custom LaTeX preamble
% PDF Standard Fonts
\usepackage{mathptmx} % Times
\usepackage[scaled=.90]{helvet}
\usepackage{courier}

%%% User specified packages and stylesheets

%%% Fallback definitions for Docutils-specific commands

% class handling for environments (block-level elements)
% \begin{DUclass}{spam} tries \DUCLASSspam and
% \end{DUclass}{spam} tries \endDUCLASSspam
\ifdefined\DUclass
\else % poor man's "provideenvironment"
  \newenvironment{DUclass}[1]%
    {% "#1" does not work in end-part of environment.
     \def\DocutilsClassFunctionName{DUCLASS#1}
     \csname \DocutilsClassFunctionName \endcsname}%
    {\csname end\DocutilsClassFunctionName \endcsname}%
\fi

% admonition environment (specially marked topic)
\ifdefined\DUadmonition
\else % poor man's "provideenvironment"
  \newbox{\DUadmonitionbox}
  \newenvironment{DUadmonition}%
    {\begin{center}
       \begin{lrbox}{\DUadmonitionbox}
         \begin{minipage}{0.9\linewidth}
    }%
    {    \end{minipage}
       \end{lrbox}
       \fbox{\usebox{\DUadmonitionbox}}
     \end{center}
    }
\fi

% field list environment (for separate configuration of `field lists`)
\ifdefined\DUfieldlist
\else
  \newenvironment{DUfieldlist}%
    {\quote\description}
    {\enddescription\endquote}
\fi

% numbered or symbol footnotes with hyperlinks and backlinks
\providecommand*{\DUfootnotemark}[3]{%
  \raisebox{1em}{\hypertarget{#1}{}}%
  \hyperref[#2]{\textsuperscript{#3}}%
}
\providecommand{\DUfootnotetext}[4]{%
  \begingroup%
  \renewcommand{\thefootnote}{%
    \protect\phantomsection\protect\label{#1}
    \protect\hyperlink{#2}{#3}}%
  \footnotetext{#4}%
  \endgroup%
}

% title for topics, admonitions, unsupported section levels, and sidebar
\providecommand*{\DUtitle}[1]{%
  \smallskip\noindent\textbf{#1}\smallskip}

% titlereference standard role
\providecommand*{\DUroletitlereference}[1]{\textsl{#1}}

% hyperlinks:
\ifdefined\hypersetup
\else
  \usepackage[hyperfootnotes=false,
              colorlinks=true,linkcolor=blue,urlcolor=blue]{hyperref}
  \usepackage{bookmark}
  \urlstyle{same} % normal text font (alternatives: tt, rm, sf)
\fi
\hypersetup{
  pdftitle={Test footnote and citation rendering},
}

%%% Body
\begin{document}
\title{Test footnote and citation rendering}
\author{}
\date{}
\maketitle

Paragraphs may contain footnote references (manually numbered
\raisebox{1em}{\hypertarget{footnote-reference-1}{}}\hyperlink{footnote-1}{[1]}, anonymous auto-numbered \raisebox{1em}{\hypertarget{footnote-reference-2}{}}\hyperlink{footnote-2}{[3]}, labeled auto-numbered \raisebox{1em}{\hypertarget{footnote-reference-3}{}}\hyperlink{label}{[2]}, or
symbolic \raisebox{1em}{\hypertarget{footnote-reference-4}{}}\hyperlink{footnote-3}{[*]}) or citation references (\cite{CIT2002}, \cite{DU2015}).
%
\DUfootnotetext{footnote-1}{footnote-reference-1}{{[}1{]}}{%
A footnote contains body elements, consistently indented by at
least 1 space.

This is the footnote's second paragraph.
}
%
\DUfootnotetext{label}{footnote-reference-3}{{[}2{]}}{%
Footnotes may be numbered, either manually (as in \raisebox{1em}{\hypertarget{footnote-reference-5}{}}\hyperlink{footnote-1}{[1]}) or
automatically using a \textquotedbl{}\#\textquotedbl{}-prefixed label.  This footnote has a
label so it can be referred to from multiple places, both as a
footnote reference (\raisebox{1em}{\hypertarget{footnote-reference-6}{}}\hyperlink{label}{[2]}) and as a \hyperref[label]{hyperlink reference}.
}
%
\DUfootnotetext{footnote-2}{footnote-reference-2}{{[}3{]}}{%
This footnote is numbered automatically and anonymously using a
label of \textquotedbl{}\#\textquotedbl{} only.

This is the second paragraph.

And this is the third paragraph.
}
%
\DUfootnotetext{footnote-3}{footnote-reference-4}{{[}*{]}}{%
Footnotes may also use symbols, specified with a \textquotedbl{}*\textquotedbl{} label.
Here's a reference to the next footnote: \raisebox{1em}{\hypertarget{footnote-reference-7}{}}\hyperlink{footnote-4}{[†]}.
}
%
\DUfootnotetext{footnote-4}{footnote-reference-7}{{[}†{]}}{%
This footnote shows the next symbol in the sequence.
}
%
\DUfootnotetext{footnote-5}{footnote-5}{{[}4{]}}{%
Here's an unreferenced footnote, with a reference to a
nonexistent footnote: \raisebox{1em}{\hypertarget{footnote-reference-8}{}}\hyperlink{footnote-6}{[5]}.
}


\section{Citations}

Here's a reference to the above, \cite{CIT2002}.
%
\DUfootnotetext{footnote-6}{footnote-reference-8}{{[}5{]}}{%
this footnote is missing in the standard example document.
}

Footnotes may contain block elements like lists  \raisebox{1em}{\hypertarget{footnote-reference-9}{}}\hyperlink{footnote-7}{[6]} \raisebox{1em}{\hypertarget{footnote-reference-10}{}}\hyperlink{list-note}{[7]} \raisebox{1em}{\hypertarget{footnote-reference-11}{}}\hyperlink{footnote-8}{[8]},
admonitions \raisebox{1em}{\hypertarget{footnote-reference-12}{}}\hyperlink{footnote-9}{[9]}, or tables \raisebox{1em}{\hypertarget{footnote-reference-13}{}}\hyperlink{footnote-10}{[10]}.
%
\DUfootnotetext{footnote-7}{footnote-reference-9}{{[}6{]}}{
\begin{enumerate}
\item An ordered list

\item in a footnote.
\end{enumerate}
}
%
\DUfootnotetext{list-note}{footnote-reference-10}{{[}7{]}}{
\begin{itemize}
\item An unordered list (bullet list)

\item in a footnote.
\end{itemize}

And a trailing paragraph.
}
%
\DUfootnotetext{footnote-8}{footnote-reference-11}{{[}8{]}}{
\begin{DUfieldlist}
\item[{Field:}]
list

\item[{with:}]
2 items.
\end{DUfieldlist}
}
%
\DUfootnotetext{footnote-9}{footnote-reference-12}{{[}9{]}}{
\begin{DUclass}{note}
\begin{DUadmonition}
\DUtitle{Note}

This is a note in a note.
\end{DUadmonition}
\end{DUclass}
\smallskip}
%
\DUfootnotetext{footnote-10}{footnote-reference-13}{{[}10{]}}{
\setlength{\DUtablewidth}{\dimexpr\linewidth-3\arrayrulewidth\relax}%
\begin{longtable*}{|p{\DUcolumnwidth{0.150}}|p{\DUcolumnwidth{0.250}}|}
\hline

a
 & 
table
 \\
\hline

in a
 & 
footnote
 \\
\hline
\end{longtable*}
}

This \raisebox{1em}{\hypertarget{footnote-reference-14}{}}\hyperlink{list-note}{[7]} is a second reference to the footnote containing
a bullet. list.

\begin{thebibliography}{CIT2002}
\bibitem[CIT2002]{CIT2002}{
Citations are text-labeled footnotes. They may be
rendered separately and differently from footnotes.
}
\bibitem[DU2015]{DU2015}{
\DUroletitlereference{Example document}, Hometown: 2015.
}
\end{thebibliography}

\end{document}
